\UnnumberedSection{Executive Summary}
\Guidance{Aim for approximately three-quarters of a page. Summarize the project, the contents of this proposal, two or three main milestones, and the estimated total cost. Write this section after drafting the body.}

AI coding agents increasingly follow \emph{agent skills}: \texttt{SKILL.md} files of natural-language instructions that an agent loads when it judges that a task matches the skill's description. Within nine months of the format's release, public GitHub repositories held 3.8 million skill files~\cite{gitskills2027}. Skills are written mainly in prose, no compiler checks them, and they spread by being copied between repositories. Software engineering has long used metrics such as size, churn, age, clone coverage, and readability to judge source code, but it is not known whether these metrics carry over to skills. Using the GitSkills dataset from the \ac{MSR} 2027 Mining Challenge~\cite{msr2027MiningChallenge}, this project asks which of these metrics have meaningful equivalents for skills and how their distributions compare with those of source code. The end product is a reproducible analysis pipeline, documented metric definitions, visualizations, and a report interpreting the findings for researchers and developers who write and reuse skills.

The pipeline loads the GitSkills tables, filters out files that are not genuine skills, runs a separate calculation script for each metric, and joins the results into a single metrics file. Each metric is adapted to skills. Size is counted in tokens, separately for the description, the body, and bundled files. Churn is split by type, because skills are more often copied and modified than edited in place. Age is measured from the first commit. Clone coverage is computed from identical content hashes. Readability scores each part of a skill with an instrument suited to it and is validated against human ratings. Performance targets include a full-dataset run within 24~h using at most 16~GB of memory, identical results on re-running, and at least 0.90 precision for cross-copy lineage links. The main milestones are a prepared and checked dataset (M1, 23~October~2026), all supported metrics calculated with initial plots (M3, 4~December~2026), and a validated, reproducible final pipeline (M6, 12~March~2027).

Every metric is first prototyped on the official 13,000-skill GitSkills sample before being scaled to the full dataset. Pilots have already run: they show that about two-thirds of skills are never edited after being added, and that skill bodies read at a median Flesch--Kincaid grade of 8.6. Testing combines unit tests against hand-calculated examples, integration tests across pipeline stages, reproducibility checks, and manual review of sample outputs. The dataset is openly licensed, all tools are open source, and computation fits on a single workstation, so the estimated total cost is \TemplateField{\$0}~CAD. The team is ready to begin. Further learning is needed in large-scale data handling, text-similarity and readability measurement, and non-parametric statistics.

\clearpage
